
\subsection*{E4.Curve: schematicity for every smooth proper curve family}

\paragraph{Theorem E4.Curve.}
Let $C\to S$ be a proper smooth finitely presented morphism of schemes
with geometrically connected one-dimensional fibres, where $S$ is an
arbitrary scheme. Then the full fppf Picard functor $\operatorname{Pic}_{C/S}$
is represented by a smooth separated $S$-scheme. This settles the
geometrically normal curve case of the last question at 57:315. It does
not settle higher-dimensional families, or geometrically integral
singular curves without further hypotheses.

\paragraph{Proof.}
The arithmetic genus is locally constant, so we work on a locus where
it is a fixed integer $g$. If $g=0$, the dual of the relative canonical
bundle is relatively ample. If $g\geq2$, the relative canonical bundle
is relatively ample. In either case the family is locally projective
on $S$, and the projective Picard theorem gives a smooth separated
Picard scheme. These are exactly the arguments of
\cite[\S9.3, Theorem~1 and Remark~2, p.~252]{BCEBLR1990}.

Suppose $g=1$. Put $E=\operatorname{Pic}^0_{C/S}$; this is an abelian
scheme by \cite[\S9.4, Proposition~4, p.~260]{BCEBLR1990}.
The degree map decomposes the Picard algebraic space as
\[
 \operatorname{Pic}_{C/S}=\coprod_{n\in\mathbb Z}\operatorname{Pic}^n_{C/S}.
\]
The diagonal divisor defines $C\to\operatorname{Pic}^1_{C/S}$, and
this map is an isomorphism: after an etale cover giving a section of
$C/S$, it is the usual genus-one Abel map. Thus $C$ is an $E$-torsor.
For $n\neq0$, the map
\[
 C\longrightarrow\operatorname{Pic}^n_{C/S},\qquad
        x\longmapsto[\mathcal O_C(nx)]                  \tag{E4.Curve.1}
\]
is equivariant for multiplication by $n$ on $E$. After an etale
cover trivializing $C$, it is a translate of $[n]:E\to E$.
Consequently it is a finite locally free $E[n]$-torsor, and
\[
        \operatorname{Pic}^n_{C/S}=C/E[n]               \tag{E4.Curve.2}
\]
as fppf sheaves. This conclusion also holds when $n$ is divisible by
a residue characteristic; etaleness of $E[n]$ is unnecessary.

We show that (E4.Curve.2) is a scheme. First suppose $S$ is locally
noetherian. The assertion is Zariski local on $S$. Fix $s\in S$ and
choose an affine open $V\subset C$ meeting $C_s$. The image of $V$
in $S$ is open because $C/S$ is smooth. Replace $S$ by an affine open
neighborhood of $s$ in that image, and replace $V$ by its inverse
image. The latter is a quasi-affine scheme. Since $C$ is an $E$-torsor
and $V$ meets every fibre, the image of $E\times_S V\to C$ is all of
$C$. Apply \cite[\S6.6, Theorem~2(c), pp.~164--166]{BCEBLR1990}:
every finite set of points of $C$ lies in one affine open. In
particular, each finite $E[n]$-orbit does. The finite locally free
quotient theorem, \cite[Proposition~66.14.1, tag~07S6]{BCEStacksFiniteQuotient},
therefore represents $C/E[n]$ by a scheme. The degree-zero piece is
$E$ itself; hence every piece and their disjoint union are schemes.
Smoothness and separatedness follow etale locally from the projective
curve theorem, or from the torsor description.

Finally, the arbitrary-base assertion follows by approximation,
without assuming that schematicity descends along arbitrary fpqc
covers. Locally take $S$ affine. The proper smooth finitely presented
scheme $C/S$ descends to a proper smooth finitely presented scheme
$C_0/S_0$ with $S_0$ noetherian. Geometric connectedness of proper
smooth fibres is an open condition and the arithmetic genus is
locally constant. Enlarging the model and restricting to the
corresponding open locus therefore arranges geometrically connected
genus-one fibres. The noetherian argument makes
$\operatorname{Pic}_{C_0/S_0}$ a scheme; its base change to $S$ represents
$\operatorname{Pic}_{C/S}$ by the definition of the relative fppf
Picard functor. This completes the proof.

\paragraph{Provenance.}
The smooth-curve argument above was already given by K\k{e}stutis
\v{C}esnavi\v{c}ius in his answer of 28 April 2015 to Poonen's question
\cite{BCECesnavicius2015}. We independently checked the cited BLR and
finite-quotient statements before locating that answer. No novelty is
claimed for Theorem E4.Curve.

\paragraph{Why the genus-one torsor counterexample route fails.}
An arbitrary torsor under an abelian scheme can be a nonschematic
algebraic space. Such a torsor alone does not give the requested
counterexample, whose original family must be a scheme. If a genus-one
scheme family $C$ is used, its multiples in torsor cohomology are
exactly (E4.Curve.2), and the affine-neighborhood theorem above makes
all of them schemes. Thus taking a scheme genus-one torsor and hoping
that $\operatorname{Pic}^n$ becomes a nonschematic multiple cannot work.
